\chapter{Lazy Updates}
The core of Dotter, is calculating the posterior probabilities of prefixes.
These posterior probabilities can be calculated exactly since both the prior, the likelihood, and the normalizing constant are all known.
However, the challenge is to make these calculations within milliseconds of the user's signal.

Our program maintains a trie of literal prefixes. Our goal is for our updates
to touch as few nodes of the trie as possible.

\section{Two Kinds of Updates}
The purpose of an update is to:
\begin{enumerate}
    \item Identify all nodes with a posterior probability exceeding the visibility threshold.
    \item Calculate the posterior probabilities of these nodes.
    This is equivalent to calculating the unnormalized posterior probability of these nodes
    (so long as they are using the same normalizing constant),
    since the posterior probability of the root node must always be 1,
     and thus we can identify the normalizing constant from the unnormalized probability of the root node.
\end{enumerate}

In order to identify visible nodes, we must know normalized posterior probabilities, which in turn requires knowing the normalizing constant, i.e., the unnormalized probability of the root node.
Hence before we can descend the trie to identify visible nodes, we must first find the unnormalized probability of the root node.
Thus our update procedure has two phases:
\begin{enumerate}
    \item Up Propagation: we choose a complete prefix partition for which we know the unnormalized probabilities, and push their sums to the root of the trie \texttt{root\_post\_Z}.
    \item Down Propagation: we descend the trie, visiting a child node if it exceeds the visibility threshold, \texttt{post\_Z - root\_post\_Z > visibility\_threshold}.
    To visit a node means to create its children if the node is a leaf, or else recalculate its children's statistics.
    (Note that when later gathering the visible nodes, we must descend the trie. A flat scan could cause a bug since \texttt{post\_Z} is not guaranteed to be correct a greater depths. I.e.,
    a deep descendant could have a higher \texttt{post\_Z} than its ancestors if there is unpushed likelihood along the way.)
\end{enumerate}

The trick is to choose the smallest prefix partition for which we can truly calculate the unnormalized posteriors.
The unnormalized posterior is a sum over one's children's unnormalized posteriors:
$$
P(x|D) = \sum_i P(x_i) P(D | x_i)
$$
In general if one term is reweighted, we will have to recalculate,
but there are three cases in which we can avoid recalculation and descent:
\begin{enumerate}
    \item Neither term (neither likelihood nor prior) changes.
    \item One term changes by a constant factor, $\sum_i (Ca_i)b_i = C \sum_i a_i b_i$.
    \item One set of terms is reweighted to a known sum, and the other set of terms is a constant value across all indices.
    $\sum_i a_i' b_i = (\sum_i a_i' / \sum_i a_i) \sum_i a_i b_i$.
\end{enumerate}
In order to find the smallest prefix partition for the Up Propagation step,
we consider the set of all nodes that satisfy any of these three cases and
then take the minimal elements of this set (resulting in an antichain / frontier of the tree).

\bigskip
We update the posterior probabilities in two ways:
\begin{enumerate}
    \item \textbf{Likelihood updates}: when the user makes a gesture, the likelihoods of the prefixes are updated. This is a traditional update.
    \item \textbf{Prior updates}: when the language model server returns new or better prior probabilities, we must recalculate the posteriors.
\end{enumerate}

\begin{table}[h]
\centering
\begin{tabular}{lll}
\hline
\textbf{} & \textbf{Likelihood Update} & \textbf{Prior Update} \\
\hline
Mathematical Term       & $P(D \mid X)$                         & $P(X)$ \\
Recompute Posteriors    & Yes                                   & Yes \\
Resize Boxes            & Yes                                   & Yes \\
How Many New Nodes?     & Dozens                                & Zero to a few \\
How to Set Phases       & Optimized / gradient descent          & Random (or ancestor / minimum placement) \\
Human-Initiated         & Yes                                   & No \\
Show Immediately        & Yes                                   & No --- wait for ancestor to finish \\
Branch of Tree Affected & Entire tree                           & Single branch \\
\hline
\end{tabular}
\caption{Comparison of likelihood updates and prior updates.}
\label{tab:updates}
\end{table}


\section{Likelihood Updates}
For our complete prefix partition we choose the set
of all nodes that were leaves in the visible trie,
along with all invisible nodes that have a visible sibling.
Equivalently, we are choosing the minimal elements of the set of nodes where
every node has no visible children (resulting in an antichain).

Since all of each of these nodes' children no visible children,
each of these nodes' children's likelihoods should be updated
in the same way as this node's likelihood (since they all share the same timer).
Since all childrens' likelihood are updated by the same factor,
this satisfies case 2 above.

In implementation we add the log likelihood of the user's gesture conditional on this node's timer
to the node's \texttt{post\_Z} and also add this value to the node's \texttt{unpushed\_likelihood}.
If later we determine that this node meets the visibility threshold and thus decide to visit it (i.e., recalculate its children's statistics),
we will add the node's \texttt{unpushed\_likelihood} to the \texttt{post\_Z} of each of its children
and reset the node's \texttt{unpushed\_likelihood} to 0.

\section{Prior Updates}
Prior updates are more subtle. To find our prefix partition
we descend from the token prefix of the update until reaching either (1) 
a node that \textit{never} was a parent to a visible node, or (2)
a space character.

If we reach a node that was never a parent to a visible node,
then we are guaranteed that its children all have the same likelihood
and thus case 3 above is satisfied.
If we reach a space character, then there is only one possible
token path up to this point. If that path does not break on the received token prefix,
then there is no change to the prior and we are in case 1 above. If that path
does contain the received token, then this node and all its literal children
will contain the received token in their paths and thus their 
priors will all change by the same factor which is case 2 above.

In implementation, we do not track the prior as a single value.
Instead each literal node tracks the prior assigned to it by each of its
(at most 15) possible token paths. The prior update id, the prior value, and the model strength
is stored for each of these paths. Unlike a single scalar \texttt{unpushed\_likelihood},
we store a set of \texttt{unpushed\_prior\_update\_ids}. When we visit a node,
we push this set to its children, and appropriately update the children's priors.
By comparing the new prior with the old prior, we can correctly modify the \texttt{post\_Z} of the children.

We want to ensure that our prior over prefixes is always a valid prior.
That means we should always use a valid prior (not the zero prior) for every next token distribution,
even if it hasn't yet been received from the language model.
A reasonable choice of default is either a uniform prior over tokens, or a 100\% prior on the terminal token.
(Note that the terminal token should \textit{never} be visited since subsequent tokens are nonsensical.)
However, this technique would also allow for the use of multiple levels of language model power (e.g. an ngram model, a provisional language model, and a full language model).

\section{Token Paths}
Since Dotter uses a letter-by-letter trie, but language models return token probabilities,
we face the task of converting a token model to a literal model.

A given literal prefix corresponds to several token paths.
For example, the prefix \texttt{sci} could be part of the words with the following
tokenizations (using the GPT-2 tokenizer):
\begin{itemize}
    \item \texttt{scientific} (single token),
    \item \texttt{science} (single token),
    \item \texttt{sci-ences} (two tokens),
    \item \texttt{sc-iss-or} (three tokens),
\end{itemize}
We can see that a literal prefix can correspond to multiple 
token prefixes (empty string, \texttt{sci}, and \texttt{sc}), and that for a given token prefix, the literal prefix
can be generated from multiple terminal tokens (\texttt{science}, \texttt{scientific}).

Formally, to convert a token model to a literal model, we take the sum across 
all compatible token prefixes $B$, and all compatible terminal tokens $T$,
$$
P(x) = \sum_{B < x} P(B) \sum_{T,\, x \leq B+T} P(T \mid B)
$$

We use two tricks to efficiently compute this sum:
\begin{enumerate}
    \item We recognize that the eligible longest token prefixes correspond to the prefixes of the last word, since tokens never contain spaces in the middle of them.
    Additionally we can limit the values of $B$ to at most 15 ancestor prefixes, since tokens are at most 15 characters long in our tokenizer.
    \item We can preconvert the language model's output of logits to an array of \textit{cumulative} probabilities.
    This allows us to efficiently compute the alphabetical sum over compatible $T$, by taking the difference between the cumulative probabilities of the start and end.
\end{enumerate}



